% Paper 3, §5 — Process reward models: the verifier leg
% Anchors: kb/notes/reasoning/process-supervision.md §1-3
%          kb/excerpts/{lightman2023-prm800k,thinkprm2025}.md
%          (wang2024-mathshepherd, she2025-r-prm, cao2025-edu-prm,
%           zheng2025-prm-survey have stub excerpts only — flagged
%           with \deepencite where load-bearing.)

\section{Process reward models: the verifier leg}
\label{sec:process-reward-models}

The third corner of the reasoning triangle is the \emph{verifier}: a
learned signal that scores a candidate solution at granularity finer
than the final answer, and uses that score to re-rank, re-weight, or
guide further computation. Its canonical artefact is the
\textbf{process \glsterm{reward-model}{reward model} (\glsterm{process-reward-model}{PRM})} — a function that scores each step
of a \glsterm{chain-of-thought}{chain-of-thought} trace
\citep[\S 2.2]{lightman2023-prm800k}.\footnote{KB excerpt:
\texttt{kb/excerpts/lightman2023-prm800k.md\#sec-2-2}.} PRMs were
introduced by \citet{lightman2023-prm800k} as the contrast to outcome
reward models (ORMs), which score only the final answer; the headline
finding of that paper — process supervision substantially outperforms
outcome supervision on the MATH benchmark — supplies the empirical
motivation for the entire 2024--2026 \glsterm{process-reward-model}{PRM} lineage. This section
formalises the per-step reward, walks the four canonical aggregation
rules, transcribes the \glsterm{prm800k}{PRM800K} experiment, traces the lineage through
Math-Shepherd's automated labelling and \glsterm{generative-prm}{ThinkPRM}'s generative-verifier
shift, and closes on the open question of whether external PRMs still
add \glsterm{value}{value} once the policy itself has been RL-trained for verifiable
reasoning.

\subsection{What a process reward model commits to}
\label{sec:prm-commitment}

Let $x$ denote a problem and $z = (s_1, s_2, \ldots, s_T)$ a reasoning
trace decomposed into $T$ steps, with each $s_t$ a contiguous span of
trace tokens (a sentence, an equation line, a paragraph). A
\textbf{process \glsterm{reward-model}{reward model}} is a function%
\footnote{KB note: \texttt{kb/notes/reasoning/process-supervision.md\#1.2}.}
\begin{equation}
  R^P_\phi : \bigl(x,\; s_1, \ldots, s_t\bigr) \;\longmapsto\; r_t \in [0, 1],
  \label{eq:prm-signature}
\end{equation}
trained to predict the per-step correctness label $c_t \in \{0, 1\}$.
The training objective is binary cross-entropy on the prefix
\citep[\S 2.2]{lightman2023-prm800k}:
\begin{equation}
  \mathcal{L}_\text{PRM} \;=\; -\;\E_{(x,\, s_{1:t},\, c_t) \sim \mathcal{D}}\!\!\left[\, c_t \log R^P_\phi(x, s_{1:t}) + (1 - c_t)\log\!\bigl(1 - R^P_\phi(x, s_{1:t})\bigr)\right].
  \label{eq:prm-loss}
\end{equation}
The conditioning is the only structural change from the \glsterm{outcome-reward-model}{ORM} loss
\citep[\S 2.1]{lightman2023-prm800k}:
\begin{equation}
  \mathcal{L}_\text{ORM} \;=\; -\;\E_{(x,\, z,\, y,\, c) \sim \mathcal{D}}\!\!\left[\, c \log R^O_\phi(x, z, y) + (1 - c)\log\!\bigl(1 - R^O_\phi(x, z, y)\bigr)\right],
  \label{eq:orm-loss}
\end{equation}
where $c \in \{0,1\}$ is the outcome correctness label on the final
answer $y$. A \glsterm{process-reward-model}{PRM} emits one score per prefix; an \glsterm{outcome-reward-model}{ORM} emits one score
per trace.

To compare a \glsterm{process-reward-model}{PRM} and an \glsterm{outcome-reward-model}{ORM} at fixed $N$ in best-of-$N$ reranking, the
$T$-vector $(r_1, \ldots, r_T)$ must be aggregated to a scalar.
Three rules dominate the literature
\citep[\S 3.1]{lightman2023-prm800k}:
\begin{align}
  R^P_{\min}(z)  \;&=\; \min_{1 \le t \le T} R^P_\phi(x, s_{1:t}), \label{eq:prm-min} \\
  R^P_{\prod}(z) \;&=\; \prod_{t=1}^{T} R^P_\phi(x, s_{1:t}), \label{eq:prm-prod} \\
  R^P_{\text{last}}(z) \;&=\; R^P_\phi(x, s_{1:T}). \label{eq:prm-last}
\end{align}
The min reading treats one bad step as fatal; the product reading
treats steps as conditionally-independent multiplicands; the last-step
reading scores only the final state. \citet[\S 4.2]{lightman2023-prm800k}
ablate all three on \glsterm{prm800k}{PRM800K}-grade verifiers and find
\textbf{min-aggregation} the strongest;\footnote{KB excerpt:
\texttt{kb/excerpts/lightman2023-prm800k.md\#sec-4-2}.} we return to
this finding in \cref{sec:prm-prm800k}.

\begin{intuition}
A \glsterm{process-reward-model}{PRM} is a per-step graded checkmark; an \glsterm{outcome-reward-model}{ORM} is one final
thumbs-up-or-down at the end of the trace. Returning to math: the
supervisory signal in \cref{eq:prm-loss} is concentrated at $T$
distinct prefixes $s_{1:t}$ for $t = 1, \ldots, T$, while
\cref{eq:orm-loss} concentrates it at a single $(z, y)$ pair. With
dense per-step labels the credit-assignment problem is solved
locally — the gradient identifies which step's residual was
mis-scored, rather than back-propagating one outcome label across an
entire $10^3$--$10^4$-token trace.
\end{intuition}

\subsection{PRM800K: the canonical process-supervision experiment}
\label{sec:prm-prm800k}

\citet{lightman2023-prm800k} construct \textbf{\glsterm{prm800k}{PRM800K}}, a dataset of
$\approx 800{,}000$ human step-correctness labels on GPT-4-generated
solutions to MATH problems \citep[\S 3]{lightman2023-prm800k}.\footnote{KB
excerpt: \texttt{kb/excerpts/lightman2023-prm800k.md\#sec-3}.} The
construction protocol: a generator (a GPT-4 variant) emits several
candidate solutions per problem; human annotators label each step as
positive, negative, or neutral; the dataset is released openly. The
resulting verifier is the foundational reference \glsterm{process-reward-model}{PRM} against which all
subsequent open-source efforts (Math-Shepherd, \glsterm{generative-prm}{ThinkPRM}, R-\glsterm{process-reward-model}{PRM},
EDU-\glsterm{process-reward-model}{PRM}) benchmark.

The headline experiment fixes $N$ across two reranking strategies and
measures best-of-$N$ accuracy on MATH test
\citep[\S 4]{lightman2023-prm800k}:\footnote{KB excerpt:
\texttt{kb/excerpts/lightman2023-prm800k.md\#sec-4}.}
\begin{align}
  \hat y_\text{ORM} \;&=\; y^{(n^\star)}, \quad n^\star \in \argmax_n R^O_\phi\!\bigl(x, z^{(n)}, y^{(n)}\bigr), \label{eq:bon-orm} \\
  \hat y_\text{PRM} \;&=\; y^{(n^\star)}, \quad n^\star \in \argmax_n R^P_{\min}\!\bigl(z^{(n)}\bigr). \label{eq:bon-prm}
\end{align}
Across the full sweep $N \in [1, 1024]$ the \glsterm{process-reward-model}{PRM} dominates the \glsterm{outcome-reward-model}{ORM} at
matched $N$, with the gap widest at moderate $N$
\citep[Fig. 4]{lightman2023-prm800k}. Both reranking strategies in
turn dominate plain \glsterm{self-consistency-2}{self-consistency}
(\cref{eq:self-consistency}) at the same sample budget. The
process-supervised verifier solves \textbf{78\%} of problems on a
representative MATH test subset
\citep[abstract]{lightman2023-prm800k}.\footnote{KB excerpt:
\texttt{kb/excerpts/lightman2023-prm800k.md\#headline}.}

The min-aggregation finding deserves emphasis. Among
\cref{eq:prm-min,eq:prm-prod,eq:prm-last}, only the min explicitly
penalises a single low-scored step regardless of context. The product
formulation, despite being the obvious likelihood-style aggregator
under a step-independence assumption, allows many high-scored steps
to mask one low-scored step in the geometric mean. The last-step rule
discards intermediate signal entirely. Empirically
\citep[\S 4.2]{lightman2023-prm800k} the min wins by a margin large
enough that it is the default aggregator in every downstream \glsterm{process-reward-model}{PRM}
paper. \citet[\S 4.3]{lightman2023-prm800k} additionally show that
\glsterm{process-reward-model}{PRM}-weighted \glsterm{self-consistency-2}{self-consistency} — majority vote with each ballot
weighted by the trace's min-\glsterm{process-reward-model}{PRM} score — improves over plain
majority vote, recovering some of the \glsterm{process-reward-model}{PRM} gain at lower verifier-cost
overhead.

\begin{intuition}
Why min wins: if a single algebraic slip turns ``$3x = 12$'' into
``$x = 5$'', every subsequent step inherits the error. The trace's
truth-\glsterm{value}{value} is bounded above by its weakest step. Mean and product
let many strong steps mask one critical error; the min surfaces it.
The empirical advantage of \cref{eq:prm-min} is consistent with
this picture, and the structural form of the gain is a
weakest-link-of-the-chain inequality: $R^P_{\min}(z) \le R^P_\phi(x, s_{1:t})$
for every $t$, with equality at the worst step.
\end{intuition}

\subsection{Math-Shepherd: removing the human-annotation bottleneck}
\label{sec:prm-mathshepherd}

The \glsterm{prm800k}{PRM800K} protocol is human-expensive --- $\approx 800$K step-labels
at expert-level math annotation cost. \citet{wang2024-mathshepherd}
remove the bottleneck by generating per-step labels automatically from
\emph{Monte Carlo rollout completion rate}.\deepencite{wang2024-mathshepherd
\S3 (excerpt skeleton-only as of 2026-05; PDF acquisition pending)} The
automatic-labelling rule, drawn from the
process-supervision KB note's lineage table%
\footnote{KB note: \texttt{kb/notes/reasoning/process-supervision.md\#3.1}.}
and the surveyed framing in \citet{zheng2025-prm-survey}, is: for each
prefix $s_{1:t}$ of a candidate trace, sample $K$ continuation rollouts
$z^{(k)} \sim p_\theta(\cdot \mid x, s_{1:t})$ to completion, extract
the answer $y^{(k)}$ from each, and define
\begin{equation}
  \tilde c_t \;=\; \frac{1}{K} \sum_{k=1}^{K} \mathbb{1}\!\bigl[y^{(k)} = y^\star(x)\bigr],
  \label{eq:mathshepherd-label}
\end{equation}
i.e., the empirical \emph{completion-success rate} from prefix
$s_{1:t}$. The label is then thresholded (or used directly as a soft
label) and substituted for the human $c_t$ in \cref{eq:prm-loss}. The
rationale is the credit-assignment heuristic: prefixes that lead to
correct answers in expectation are good prefixes; prefixes that do
not are bad. Outcome supervision on the rollouts implicitly produces
process supervision on the prefixes.

The implication is structural: any reasoning policy plus a verifiable
outcome checker (an arithmetic verifier, a Python interpreter, a
regex on the boxed final answer) can manufacture its own \glsterm{process-reward-model}{PRM} training
data. The human-annotation cost collapses from $O(\text{PRM800K})$
to the compute cost of $K \times T$ rollout extensions per training
example. Math-Shepherd-style automatic labelling is the substrate
that subsequent generative-\glsterm{process-reward-model}{PRM} efforts (\cref{sec:prm-thinkprm}) and
reasoning-trained PRMs (\cref{sec:prm-r-edu}) build on.

\begin{analogy}
Math-Shepherd is the \glsterm{process-reward-model}{PRM} equivalent of Q-learning's bootstrapped
\glsterm{value}{value} estimate: the \glsterm{value}{value} of a state is approximated by the
empirical return-from-state under the current policy, rather than by
human labelling of the state. Returning to math: \cref{eq:mathshepherd-label}
is a Monte Carlo estimator of the true success probability
$\Pr_{z \sim p_\theta(\cdot \mid x, s_{1:t})}[y^\star(x)]$, and the
\glsterm{process-reward-model}{PRM} trained on \cref{eq:prm-loss} with $\tilde c_t$ in place of $c_t$
is a learned approximator to that conditional success probability.
The bias-variance characteristics are the standard ones of $K$-sample
Monte Carlo.
\end{analogy}

\subsection{ThinkPRM: process reward models that think}
\label{sec:prm-thinkprm}

The 2025 standout shift: instead of a classifier head that emits one
scalar per prefix, the verifier is itself a small reasoning \glsterm{llm}{LLM} that
reads the prefix, emits its own short verification \glsterm{chain-of-thought}{chain-of-thought},
and outputs a verdict token \citep[\S 3]{thinkprm2025}.\footnote{KB
excerpt: \texttt{kb/excerpts/thinkprm2025.md\#sec-3}.} Formally,
\glsterm{generative-prm}{ThinkPRM} replaces the discriminative $R^P_\phi$ of \cref{eq:prm-signature}
with a generative verifier
\begin{equation}
  R^{P,\text{gen}}_\phi(x, s_{1:t}) \;=\; p_\phi\!\bigl(\langle\text{verdict}\rangle = \mathsf{correct} \;\big|\; x,\; s_{1:t},\; v\bigr), \quad v \sim p_\phi\!\bigl(\cdot \mid x, s_{1:t}\bigr),
  \label{eq:thinkprm-gen}
\end{equation}
where $v$ is the verifier's own reasoning trace over the candidate
prefix.

The headline result: a long-\glsterm{chain-of-thought}{CoT} verifier fine-tuned on
\textbf{1\%} of \glsterm{prm800k}{PRM800K} (roughly 8K step-labels rather than 800K)
plus self-generated synthetic verification traces beats discriminative
verifiers trained on the full \glsterm{prm800k}{PRM800K}
\citep[abstract]{thinkprm2025}.\footnote{KB excerpt:
\texttt{kb/excerpts/thinkprm2025.md\#abstract}.} On out-of-domain
evaluation, \glsterm{generative-prm}{ThinkPRM} exceeds the full-\glsterm{prm800k}{PRM800K} discriminative baseline
by \textbf{8\%} on a GPQA-Diamond subset and \textbf{4.5\%} on
LiveCodeBench, and beats \glsterm{llm}{LLM}-as-a-Judge by \textbf{7.2\%} on
\glsterm{processbench}{ProcessBench} at matched verification-token budget
\citep[abstract]{thinkprm2025}.\footnote{KB excerpt:
\texttt{kb/excerpts/thinkprm2025.md\#headline}.} The data-efficiency
ratio is two orders of magnitude on labels, traded for verification
compute.

The cost asymmetry is real: best-of-$N$ with a \glsterm{discriminative-prm}{discriminative PRM}
costs $O(N)$ classifier passes; with a generative \glsterm{process-reward-model}{PRM},
$O(N \times L_\text{verify})$ tokens, where $L_\text{verify}$ is the
verification trace length. The asymmetry is the same one that
distinguishes a \glsterm{chain-of-thought}{chain-of-thought} policy from a direct-answer policy
(\cref{sec:cot-phenomenon}), applied now to the verifier role: spend
inference compute inside the verifier to extract more signal per
labelled example. The mechanism: a verifier's \glsterm{chain-of-thought}{CoT} can apply
algebraic-verification, type-checking, and dimensional-analysis
checks that a feed-forward classifier cannot internalise from
labelled examples alone \citep[\S 5]{thinkprm2025}.\footnote{KB
excerpt: \texttt{kb/excerpts/thinkprm2025.md\#sec-5}.}

\subsection{R-PRM and EDU-PRM: the 2025 PRM lineage}
\label{sec:prm-r-edu}

The 2025 \glsterm{process-reward-model}{PRM} literature splits along two axes orthogonal to the
discriminative-vs-generative split:

\paragraph{Reasoning-aware PRMs.} \citet{she2025-r-prm} train a \glsterm{process-reward-model}{PRM}
that emits a short reasoning step before scoring; the architecture is
discriminative (a verdict head) but the input now includes a brief
verifier-reasoning prefix. This generalises better across distribution
shifts than vanilla discriminative PRMs.\deepencite{she2025-r-prm \S3
(excerpt skeleton-only as of 2026-05)} The contrast with
\cref{eq:thinkprm-gen} is one of degree: \glsterm{generative-prm}{ThinkPRM}'s verifier emits a
full long-\glsterm{chain-of-thought}{CoT} trace and a verdict token; R-\glsterm{process-reward-model}{PRM} emits a short
reasoning preamble and a discriminative score. The two converge as
the preamble length grows.

\paragraph{Entropy-driven step boundaries.} \citet{cao2025-edu-prm}
challenge the common assumption that steps are fixed by the policy's
linebreaks or paragraph markers. EDU-\glsterm{process-reward-model}{PRM} defines step boundaries by a
local entropy spike in the policy's per-token distribution: a step
ends where the model is most uncertain about the next token, on the
intuition that high-entropy positions correlate with reasoning
branch-points. Trained on \textbf{1.5\%} of the data of Math-Shepherd,
EDU-\glsterm{process-reward-model}{PRM} matches or beats Math-Shepherd on
\glsterm{processbench}{ProcessBench}.\deepencite{cao2025-edu-prm \S3 (excerpt skeleton-only as
of 2026-05)} The framing recasts step labelling from a syntactic
problem (where do paragraphs end?) to an information-theoretic one
(where does the model branch?).

The two strands jointly suggest that the 2023 \glsterm{prm800k}{PRM800K} formulation —
human-labelled, syntactically-bounded, classifier-head — was a
specific, easily-relaxable instantiation of a more general per-step
verification programme. Each relaxation (Math-Shepherd, \glsterm{generative-prm}{ThinkPRM},
R-\glsterm{process-reward-model}{PRM}, EDU-\glsterm{process-reward-model}{PRM}) trades human labels or syntactic priors for some
combination of automated rollouts, verifier compute, and
information-theoretic boundaries.

\subsection{The 2025 PRM survey landscape}
\label{sec:prm-survey}

\citet{zheng2025-prm-survey} catalogue the 2023--2025 \glsterm{process-reward-model}{PRM} landscape
along the discriminative-vs-generative axis and the
test-time-vs-training-time axis.\deepencite{zheng2025-prm-survey
overview (excerpt skeleton-only as of 2026-05)} The survey's
synthesised finding (drawn from the KB note's
treatment%
\footnote{KB note: \texttt{kb/notes/reasoning/process-supervision.md\#3.2}.})
runs in two parts. First, generative PRMs win on competition-math
benchmarks where each step has a verifiable certificate
(algebraic verification, equation manipulation, type-checking) the
verifier's \glsterm{chain-of-thought}{CoT} can apply; discriminative PRMs remain competitive on
broader domains where step-level checks are not formal. Second, and
more consequentially, the survey reports that most PRMs improve
test-time-compute performance \emph{until the base model becomes a
strong \glsterm{reasoning-model}{reasoning model} itself} — at which point the marginal gain
from an external \glsterm{process-reward-model}{PRM} at test time decreases. This is the opening of
the contradiction discussed in \cref{sec:prm-post-r1}.

\subsection{PRM utility post-R1: the open contradiction}
\label{sec:prm-post-r1}

\citet{deepseek-r1} achieved the headline AIME 2024 jump from
$15.6\%$ to $77.9\%$ pass\@1 with \emph{terminal verifier rewards
only} — no \glsterm{process-reward-model}{PRM} in the RL loop \citep[\S 2.2]{deepseek-r1}.\footnote{KB
excerpt: \texttt{kb/excerpts/deepseek-r1.md\#sec-2-2-reward}.} The
R1 reward is rule-based: an answer-match check on math problems and a
compiler / unit-test check on code, returning a single binary outcome
reward per trajectory. PRMs do not appear in the headline RL loop of
the strongest open-weight \glsterm{reasoning-model}{reasoning model}.

The implication, surfaced jointly by \citet{thinkprm2025} and
\citet{zheng2025-prm-survey}: as the policy itself internalises
verification through long-\glsterm{chain-of-thought}{CoT} RL, the marginal \glsterm{value}{value} of an external
\glsterm{process-reward-model}{PRM} at test time shrinks. The 2026 picture is mixed; the contradiction
deserves explicit treatment.

\begin{contradiction}
\textbf{\glsterm{process-reward-model}{PRM} utility post-R1 — diminishing or persistent?}
The \glsterm{prm800k}{PRM800K}-era headline ``\glsterm{process-reward-model}{PRM} > \glsterm{outcome-reward-model}{ORM} at matched $N$'' was established
on a non-reasoning-trained generator (a GPT-4 variant)
\citep[\S 4]{lightman2023-prm800k}. As the policy class moves to
\glsterm{rlvr-3}{RLVR}-trained reasoning models, the picture shifts. Some 2026
measurements show \glsterm{process-reward-model}{PRM}-best-of-$N$ over R1-distill outputs still
delivers gains; others find diminishing returns as the policy's own
internal verification approaches the \glsterm{process-reward-model}{PRM}'s signal
\citep[\S 5]{thinkprm2025}. The 2025 \glsterm{process-reward-model}{PRM} survey
\citep{zheng2025-prm-survey} reports that most PRMs improve
test-time-compute performance until the base model becomes a strong
\glsterm{reasoning-model}{reasoning model} itself. Whether external PRMs survive the
reasoning-trained-model era as a training signal --- as opposed to
inference-time rerankers --- is unresolved as of 2026-Q2; \cref{sec:prm-survival}
revisits the question with the post-R1 evidence.
\end{contradiction}

A second, related tension: \emph{process hacking}. \glsterm{process-reward-model}{PRM}-shaped \glsterm{rlvr-3}{RLVR}
\citep[\S 2.4]{lightman2023-prm800k}\footnote{KB note:
\texttt{kb/notes/reasoning/process-supervision.md\#4.2}.} replaces or
augments the binary outcome reward of \glsterm{rlvr-3}{RLVR} (\cref{sec:rlvr}) with a
per-step \glsterm{process-reward-model}{PRM} score. The middle-ground appeal is theoretical: dense
signal beats sparse signal for credit assignment. The empirical risk:
a misaligned \glsterm{process-reward-model}{PRM} rewards steps that \emph{look} correct without
producing the right answer, the step-level analogue of outcome-hacking
under a learned reward. Whether \glsterm{process-reward-model}{PRM}-shaped \glsterm{grpo-2}{GRPO} actually outperforms
verifier-only \glsterm{grpo-2}{GRPO} at matched compute is benchmark-specific and
unsettled in 2026.

\subsection{Forward link}
\label{sec:prm-forward}

The \glsterm{process-reward-model}{PRM} scores per-step \glsterm{value}{value} $r_t \in [0, 1]$ on a partial prefix
$s_{1:t}$. With this signal in hand, the inference decoder can do
more than rerank --- it can \emph{steer expansion} during generation.
\Cref{sec:tree-search} introduces the tree-search families (best-first,
beam, MCTS) that use a \glsterm{process-reward-model}{PRM} as the \glsterm{value}{value} function $V_\phi(s_{1:t}) =
R^P_\phi(x, s_{1:t})$ in the canonical AlphaZero-style recurrence,
formalises the rStar-Math UCT rule, and presents the cost analysis
of $O(N_\text{rollout} \cdot D \cdot L_\text{step})$ per simulation.
The aggregation rules of
\cref{eq:prm-min,eq:prm-prod,eq:prm-last} carry over: tree-search
nodes are scored by $R^P_{\min}$ over their root-to-node prefix, so
the weakest-link inequality of \cref{sec:prm-prm800k} continues to
shape which branches survive expansion. The \glsterm{process-reward-model}{PRM} is then re-encountered
in \cref{sec:prm-survival} as the centerpiece of the post-R1 survival
question --- whether external verifiers still earn their training and
inference cost once the policy itself reasons.
